%------------------------------------------------------------------------------%

\usepackage{integracao_e_series}
\usepackage{cancel}

\title{Séries Numéricas -- Testes de Convergência}

%------------------------------------------------------------------------------%
\begin{document}

  \addtitle

%------------------------------------------------------------------------------%
%------------------------------------------------------------------------------%
\section{Testes da Divergência}
%------------------------------------------------------------------------------%
%------------------------------------------------------------------------------%

%------------------------------------------------------------------------------%
\begin{frame}{Testes de Convergência}
  \begin{itemize}[<+->]
    \item Em geral não temos uma fórmula fechada para $S_n$ \\[2mm]
    \item Queremos saber se a série converge \\[2mm]
    \item Abrimos mão de calcular sua soma exata \\[2mm]
    \item Podemos calcular aproximações computacionalmente
  \end{itemize}
\end{frame}

%------------------------------------------------------------------------------%
\begin{frame}{Termo Geral de uma Serie Convergente}
  Como somamos infintos termos se eles não forem ``pequenos'' a série não converge

  \pause
  \vspace{\baselineskip}
  Para que \(\;\sum_{n=1}^\infty a_n\;\) seja convergente $a_n$ precisa ir para zero
\end{frame}

%------------------------------------------------------------------------------%
\begin{frame}{Limite do Termo Geral}

  Se a série \(\sum_{n=1}^\infty a_n\) é convergente, então $a_n\to 0$

  \pause
  \vspace{\baselineskip}
  Como a série converge
  \[ \lim S_n  = S \]

  \pause
  Calculando
  \begin{align*}
    \onslide<+->{\lim a_n & = \lim(S_n-S_{n-1})       \\}
    \onslide<+->{         & = \lim S_n - \lim S_{n-1} \\}
    \onslide<+->{         & = S - S = 0}
  \end{align*}
\end{frame}

%------------------------------------------------------------------------------%
\begin{frame}{Teste da Divergência}
  Se a sequência $a_n$ não converge para zero, então a série $\sum a_n$ diverge
\end{frame}

%------------------------------------------------------------------------------%
\section{Lista Mínima}

%------------------------------------------------------------------------------%
\begin{frame}{Lista Mínima}
  Estudar as Seção 6.3 da Apostila

  \vspace{\baselineskip}
  Exercícios: 3a-d

  \vfill
  Atenção: A prova é baseada no livro, não nas apresentações
\end{frame}

%------------------------------------------------------------------------------%
\end{document}
%------------------------------------------------------------------------------%
